% Generated draft: requires mathematical and source review.
\documentclass[11pt]{article}
\usepackage[margin=25mm]{geometry}
\usepackage{fontspec}
\setmainfont{Latin Modern Roman}
\usepackage{amsmath,amssymb,mathtools,unicode-math}
\setmathfont{Latin Modern Math}
\usepackage{enumitem,needspace,xurl,hyperref}
\setlength{\parindent}{0pt}
\setlength{\parskip}{5pt}
\setlength{\emergencystretch}{3em}
\newcommand{\N}{\mathbb{N}}
\newcommand{\Z}{\mathbb{Z}}
\newcommand{\Q}{\mathbb{Q}}
\newcommand{\R}{\mathbb{R}}
\newcommand{\C}{\mathbb{C}}
\newcommand{\F}{\mathbb{F}}
\newcommand{\A}{\mathbb{A}}
\newcommand{\PP}{\mathbb{P}}
\newcommand{\E}{\mathbb{E}}
\newcommand{\eps}{\varepsilon}
\newcommand{\dd}{\,\mathrm{d}}
\begin{document}
\section*{UnsolvedMath ID: 30003351}
\textbf{Source problem code:} catalogue entry
\subsection{Definitions and mathematical statement}
For a specified smooth manifold \(M\), write
\[
\operatorname{Met}(M)=\left\{g\in\Gamma^\infty(\operatorname{Sym}^2T^*M):g_x(v,v)>0\text{ for every }x\in M\text{ and }0\ne v\in T_xM\right\}
\]
for its smooth Riemannian metrics. To discuss connected components, this set, or a specified subspace \(\mathcal X\subseteq\operatorname{Met}(M)\), must be equipped with a topology. If a specified group \(G\) of diffeomorphisms preserves \(\mathcal X\), its action by pullback gives a moduli space \(\mathcal X/G\), equipped with the quotient topology. For a topological space \(Y\), let \(\operatorname{Comp}(Y)\) denote the set of its connected components. Connected components are not to be identified with path components without further justification.

The supplied research question is: \emph{How can one detect the connected components of the space of Riemannian metrics or its moduli space?} In the notation above, this asks for ways to distinguish elements of \(\operatorname{Comp}(\mathcal X)\) or \(\operatorname{Comp}(\mathcal X/G)\). An invariant constant on each connected component can distinguish two such components when it takes different values on them; an invariant on the moduli space must also be invariant under the specified group action. The source does not say whether a complete classification or only sufficient criteria for distinguishing some components are sought.

These definitions describe possible parameters of the question, not additional source hypotheses. The record does not specify \(M\), its dimension, compactness, boundary conditions, the topology on the metric space, the acting group, or any curvature or completeness restriction defining \(\mathcal X\). Its background discusses curvature-restricted moduli spaces, whereas its statement mentions Riemannian metrics without a restriction. Thus the intended domain and quantifiers remain unresolved, and the record does not determine a single precise conjecture.
\subsection{Short English statement}
The question asks for methods that distinguish connected components of spaces of Riemannian metrics and their moduli spaces. A precise target requires specifying the manifolds, allowed metrics, topology, and equivalence relation. The supplied background suggests curvature restrictions, but the statement does not impose them. This is a statement-only draft requiring human mathematical review.
\subsection{Sources}
\begin{itemize}
\item Local UnsolvedMath registry: supplied record with permanent ID \(30003351\), title \emph{Detecting Components of Riemannian Metric Moduli Spaces}, source code \texttt{catalogue entry}, in the Oberwolfach Reports Open Problems collection. No individual proposer is supplied; the recorded proposed year is 2017.
\item Supplied report citation: \emph{Mini-Workshop: Spaces and Moduli Spaces of Riemannian Metrics (2018)}, Oberwolfach Reports, pp.\ 133--166, DOI \texttt{10.4171/owr/2017/3}. \url{https://doi.org/10.4171/owr/2017/3}. The record describes these reports as documenting mathematical workshops and their open-problem discussions. The citation's 2018 label and the record's 2017 dating are retained without reconciliation.
\item Supplied literature source, labeled ``Oberwolfach report on disconnected moduli'': \url{https://doi.org/10.4171/owr/2017/3}.
\item Supplied literature source, labeled ``Nonconnected moduli spaces of nonnegative sectional curvature'': \url{https://arxiv.org/abs/1601.04877}.
\item Supplied literature source, labeled ``Kreck--Stolz nonconnected moduli spaces'': \url{https://doi.org/10.2307/2152742}.
\item Supplied literature assessment dated exactly \texttt{2026-08-21}: the record labels the problem open and states that index, eta, and \(s\)-invariants detect components in many curvature-restricted moduli spaces, while asserting that no general component-detection theory for spaces of metrics has been proved. This is supplied dated curation derived from the OWR material, not an independently verified literature assessment. The record expressly requests independent verification of time-sensitive status and cited claims.
\end{itemize}
\end{document}
